%% LyX 2.4.4 created this file.  For more info, see https://www.lyx.org/.
%% Do not edit unless you really know what you are doing.
\documentclass[english]{article}
\usepackage[T1]{fontenc}
\usepackage[latin9]{inputenc}
\usepackage{enumitem}
\usepackage{amsmath}
\usepackage{amsthm}
\usepackage{amssymb}
\usepackage{geometry}
\geometry{verbose,tmargin=1in,bmargin=1in,lmargin=1in,rmargin=1in}

\makeatletter
%%%%%%%%%%%%%%%%%%%%%%%%%%%%%% Textclass specific LaTeX commands.
\numberwithin{equation}{section}
\numberwithin{figure}{section}
\newlength{\lyxlabelwidth}      % auxiliary length 

%%%%%%%%%%%%%%%%%%%%%%%%%%%%%% User specified LaTeX commands.
\usepackage{fancyhdr}
\usepackage{lastpage}
\pagestyle{fancy}
\usepackage{enumitem}
\usepackage{ifthen}
\everymath=\expandafter{\the\everymath\displaystyle}
%\DeclareMathSizes{10}{10}{10}{10}
\usepackage{multicol}

\usepackage{tikz}
\usetikzlibrary{patterns}
\usetikzlibrary{plotmarks}
\usepackage{pgfplots}
\pgfplotsset{compat=newest}

\usepackage{calc}
\usepackage[nomessages]{fp}

\usepackage{datetime}
% Added by lyx2lyx
\setlength{\parskip}{\smallskipamount}
\setlength{\parindent}{0pt}

\makeatother

\usepackage{babel}
\begin{document}
\renewcommand{\author}{Dr.\,James Ashe}

\newcommand{\classNum}{335}

\newcommand{\classSec}{A}

\newcommand{\nameType}{Solutions}

\newcommand{\examType}{Homework 1}

\newcommand{\Date}{\formatdate{28}{1}{2014}} %%\AdvanceDate[12] \today

%\newdate{my_date}{\day}{\month}{\year}%   
%\AdvanceDate[-5] 
%\displaydate{my_date}

%%First page's header%%%%%%%%%%%%%%%%%%%%%%
\fancypagestyle{firststyle} %%header settings for first pagr
{
	\fancyhead[L]{MTH \classNum{}(\classSec{}) \author{}\\
	\textbf{\examType{}} \Date{}}
	\fancyhead[C]{\textbf{\nameType}}
	\setlength{\headsep}{14pt}
}
%%%%%%%%%%%%%%%%%%%%%%%%%%%%%%%%%%%%%%%%%%

%%Next page's header%%%%%%%%%%%%%%%%%%%%%%%
%\setlist{nolistsep} %eliminates space between list items
\lhead{MTH \classNum{}(\classSec{})} 
\chead{\textbf{\nameType}} 
\cfoot{\thepage{}~of~\pageref{LastPage}} 
\setlength{\headsep}{5pt} 
%%%%%%%%%%%%%%%%%%%%%%%%%%%%%%%%%%%%%%%%%%%

%%Various settings; see the preamble for other settings%%%%%
%%\setenumerate[0]{leftmargin=12pt,itemindent=0pt,labelwidth=0pt}
\renewcommand{\labelenumi}{\textbf{\arabic{enumi}.}}
\renewcommand{\labelenumii}{\textbf{\alph{enumii}.}}
\thispagestyle{firststyle}
%%%%%%%%%%%%%%%%%%%%%%%%%%%%%%%%%%%%%%%%%%

\textbf{Homework}
\begin{enumerate}
\item Complete the following.

\begin{enumerate}
\item Write in tabular form: 

\begin{enumerate}
\item $A=\{x\colon x\mbox{ is a letter in "do not repeat"}\}$

\begin{enumerate}
\item []$A=\{d,o,n,t,r,e,p,a\}$
\end{enumerate}
\item $3\mathbb{Z}=\left\{ x\in\mathbb{Z}\colon x=3n\mbox{ for some \ensuremath{n\in\mathbb{Z}}}\right\} $

\begin{enumerate}
\item []$3\mathbb{Z}=\{\dots,-6,-3,0,3,6,\dots\}$
\end{enumerate}
\item $N=\left\{ x\colon x\mbox{ is odd and \ensuremath{x}is even}\right\} $
(recall that $0$ is even)

\begin{enumerate}
\item []$N=\emptyset$
\end{enumerate}
\end{enumerate}
\item Write these sets in set-builder form:

\begin{enumerate}
\item Let $A=\{3,4,5,6,7,\dots\}$

\begin{enumerate}
\item []$A=\{x\in\mathbb{Z}\colon x\geq3\}$
\end{enumerate}
\item Let $B=\{3,5,7,9,\dots\}$

\begin{enumerate}
\item []$B=\{x\colon x=2n+1\mbox{ where }n\in\mathbb{Z}\mbox{ and }n\geq1\}$ 
\end{enumerate}
\item Let $\mathbb{Q}$ be the set of rational numbers.

\begin{enumerate}
\item []$\mathbb{Q}=\{x\colon x=\frac{a}{b}\mbox{ where }a,b\in\mathbb{Z}\}$
\end{enumerate}
\end{enumerate}
\end{enumerate}
\item Find all subsets of $A=\{1,2,3\}$.

\begin{enumerate}
\item []$\{1,2,3\}$, $\{1,2\}$, $\{1,3\}$, $\{2,3\}$, $\{1\}$, $\{2\}$,
$\{3\}$, $\emptyset$. There are $3$ elements in $A$ so there are
$2^{3}=8$ subsets.
\end{enumerate}
\item Prove that every nonempty subset has a proper subset. 

\begin{proof}
Let $A$ be any non-empty subset. $\emptyset\subseteq A$ since $\emptyset$
is a subset of every set, and $\emptyset\subset A$ (a proper subset)
since $A\neq\emptyset$. 
\end{proof}
\item Prove that if $A$ is a subset of $\emptyset$, then $A=\emptyset$.

\begin{proof}
Let $A\subseteq\emptyset$. $\emptyset\subseteq A$ for any subset
$A$. So then $A=\emptyset$.
\end{proof}
\item Prove that $A=\{2,3,4,5\}$ is not a subset of $B=\{x\mid x\mbox{ is even}\}$.

\begin{proof}
$5$ is not even since $5=2\cdot2+1$. So then $5\notin B$. Thus
$A\nsubseteq B$.
\end{proof}
\item Prove that $A$ and $B$ are subsets of $A\cup B$.

\begin{proof}
Let $a\in A$. $a\in A$ so then $a\in A\cup B$ by definition of
the union of sets. Thus $A\subseteq A\cup B$. Similarly we can show
$B\subseteq A\cup B$. 
\end{proof}
\item Prove that $A=A\cup A$ and $A=A\cap A$.

\begin{proof}
$\left(A=A\cup A\right)$. By exercise 6, $A\subseteq A\cup A$. So
we need only show that $A\cup A\subseteq A$. Let $a\in A\cup A$.
Then $a\in A$ or $a\in A$, implying that $A\cup A\subseteq A$.
Thus $A\cup A=A$

$\left(A=A\cap A\right)$. Let $a\in A$. Then $a\in A$ and $a\in A$.
Thus $A\subseteq A\cap A$. Let $a\in A\cap A$. Then $a\in A$ and
$a\in A$ by definition, and this implies that $A\cap A\subseteq A$. 

Thus $A=A\cap A$.
\end{proof}
\item Prove that $A\cap B$ is a subset of $A$ and of $B$.

\begin{proof}
Let $x\in A\cap B$. Then $x\in A$ and $x\in B$ by definition. So
then $x\in A$ implying that $A\subseteq A\cap B$. Similarly we can
show that $B\subseteq A\cap B$.
\end{proof}
\item Prove that $A\cup B=\emptyset$ implies that $A=\emptyset$ and $B=\emptyset$.

\begin{proof}
Suppose that $A\cup B=\emptyset$. $A\subseteq A\cup B=\emptyset$
by exercise 6. Thus $A\subseteq\emptyset$, and this implies that
$A=\emptyset$ by exercise 4.

Similarly we can show that $B=\emptyset$.
\end{proof}
\item Prove that $A\cup B=B$ if and only if $A\subseteq B$.

\begin{proof}
$(\Rightarrow)$ Let $A\cup B=B$. $A\subseteq A\cup B=B$ (by exercise
6) which implies that $A\subseteq B$.\\
$(\Leftarrow)$ Let $A\subseteq B$. $B\subseteq A\cup B$ (again
exercise 6) so we only need to show that $A\cup B\subseteq B$. Let
$a\in A\cup B$ then $a\in A$ or $a\in B$. If $a\in B$ then $A\cup B\subseteq B$
as required. If $a\in A$ then $A\subseteq B$ implies that $a\in B$.
Thus $A\cup B\subseteq B$. Which implies that $A\cup B=B$.
\end{proof}
\begin{itemize}
\item For ``if and only if/$\iff$'' statements, you need to show both
directions, $\Rightarrow$ and $\Leftarrow$. 
\item For $A=B$ statements, it is usually easiest to show that $A\subseteq B$
and $B\subseteq A$ separately.
\end{itemize}
\item Prove that $A\cap B=B$ if and only if $B\subseteq A$.

\begin{proof}
$(\Rightarrow)$ Let $A\cap B=B$. $B=A\cap B\subseteq A$ by exercise
8 so then $B\subseteq A$.\\
$(\Leftarrow)$ Let $B\subseteq A$. We already know that $A\cap B\subseteq B$
be ex.\,8. So we need only show that $B\subseteq A\cap B.$ Let $b\in B.$
By assumption $B\subseteq A$ so then $b\in A$ which implies that
$b$ is in $A$ and $B$, i.e., $b\in A\cap B$. 
\end{proof}
\item Prove that $A\cup\left(A\cap B\right)=A$.

\begin{proof}
$A\cup(A\cap B)=A$ if and only if $A\cap B\subseteq A$ by ex.\,10,
and $A\cap B\subseteq A$ by ex.\,8.
\end{proof}
\item Prove that $A\cap\left(A\cup B\right)=A$.

\begin{proof}
$A\cap(A\cup B)=A$ if and only if $A\subseteq A\cup B$ by ex.\,11,
and $A\subseteq A\cup B$ by ex.\,6.
\end{proof}
\item Prove that $\mathcal{P}(A\cap B)=\mathcal{P}(A)\cap\mathcal{P}(B)$.

\begin{proof}
Let $P\in\mathcal{P}(A\cap B)$. Then $P\subseteq A\cap B$ implying
that $P\subseteq A$ and $P\subseteq B$ thus $P\in\mathcal{P}(A)\cap\mathcal{P}(B)$.
\\
Now let $P\in\mathcal{P}(A)\cap\mathcal{P}(B)$. Thus $P\subseteq A$
and $P\subseteq B$ which means that $P\subseteq A\cap B$ and that
$P\in\mathcal{P}(A\cap B$).
\end{proof}
\end{enumerate}
\vfill

\textbf{Presentation problems.}
\begin{enumerate}
\item Prove that $A-B$ is a subset of $A\cup B$.
\item Prove that if $A$ a subset of $B$ then $A\cap B=A$.
\item Let $A\cap B=\emptyset$. Prove that $B\cap A'=B$.
\item Prove that if $A$ is a subset of $B$ then $A\cup B=B$.
\item Prove that $A'-B'=B-A$.
\item Prove that if $A$ is a subset of $B$ then $B'$ is a subset of $A'$.
\item Prove that if $A\cap B=\emptyset$ then $A\cup B'=B'$.
\item Prove that $\left(A\cap B\right)'=A'\cup B'$.
\item Prove that $A\subset B$ implies that $A\cup\left(B-A\right)=B$.
\end{enumerate}

\end{document}
